%% ma: L10-B13-0115
%% lop: 10
%% chuong: 5
%% bai: 13
%% dang: 10.13.10
%% loai: TLN
%% muc: TH
%% dap_an: "174"
%% nguon: "chuyen-de-cac-so-dac-trung-cua-mau-so-lieu-khong-ghep-nhom-toan-10-chuong-trinh-moi.pdf (GV. Phan Nhật Linh) (P03), Bài 02 Dạng 2 Phần III Câu 6"
%% kien_thuc: "Lớp 10 Bài 13 (trung vị, tứ phân vị)"
%% trang_thai: chinh-thuc
%% ghi_chu: "Cùng số liệu Phần I Câu 4 nhưng khác loại câu. Tự duyệt theo yêu cầu giáo viên 22/10/2026."
\begin{ex}
Chiều cao (cm) của một nhóm học sinh là: $160$; $178$; $150$; $164$; $168$; $176$; $156$; $172$. Tìm tứ phân vị thứ ba của mẫu số liệu.
\shortans{174}
\loigiai{Sắp xếp: $150$; $156$; $160$; $164$; $168$; $172$; $176$; $178$. Nửa trên $168;172;176;178$ nên $Q_3=\dfrac{172+176}2=174$.}
\end{ex}
